Une idée élémentaire pour découvrir la matrice \squote{magique} $A$ est de noter que nous pouvons écrire
$x^2 - 8 y^2
=
\begin{pmatrix}
  x & y
\end{pmatrix}
Q
\begin{pmatrix}
  x \\
  y
\end{pmatrix}$
en posant
$Q
=
\begin{pmatrix}
  1 & 0  \\
  0 & -8
\end{pmatrix}$.%
\footnote{
	Le lecteur connaissant les formes quadratiques ne sera pas surpris par cette réécriture.
}


\medskip

Notant
$\begin{pmatrix}
  x' \\
  y'
\end{pmatrix}
=
A
\begin{pmatrix}
  x \\
  y
\end{pmatrix}$,
nous avons
$\begin{pmatrix}
  x' & y'
\end{pmatrix}
Q
\begin{pmatrix}
  x' \\
  y'
\end{pmatrix}
=
\begin{pmatrix}
  x & y
\end{pmatrix}
A^T Q A
\begin{pmatrix}
  x \\
  y
\end{pmatrix}$.
Il devient alors naturel de chercher $A$ vérifiant $A^T Q A = Q$,
puisque d'une solution
$\begin{pmatrix}
  x \\
  y
\end{pmatrix}$
nous pourrons passer à une \squote{autre} solution
$\begin{pmatrix}
  x' \\
  y'
\end{pmatrix}
=
A \begin{pmatrix}
  x \\
  y
\end{pmatrix}$
comme dans le sujet du Bac S.


\medskip

Supposons donc avoir $A^T Q A = Q$.
Utilisant le déterminant, nous avons comme contrainte immédiate que $\det A = \pm 1$ \textit{($A$ doit donc être inversible)}.


% ------------------ %


\medskip

\textit{\textbf{Cas 1:} supposons d'abord $\det A = 1$.}

\medskip

Notant
$A = \begin{pmatrix}
  a & b \\
  c & d
\end{pmatrix}$,
de sorte que
$A^{-1} = \begin{pmatrix}
  d  & -b \\
  -c & a
\end{pmatrix}$,
nous avons:

\begin{flalign*}
	A^T Q A = Q & \Longleftrightarrow  A^T Q = Q A^{-1} & \\
	            & \Longleftrightarrow
	\begin{pmatrix}
	  a & c \\
	  b & d
	\end{pmatrix}
	\begin{pmatrix}
	  1 & 0  \\
	  0 & -8
	\end{pmatrix}
	=
	\begin{pmatrix}
	  1 & 0  \\
	  0 & -8
	\end{pmatrix}
	\begin{pmatrix}
	  d  & -b \\
	  -c & a
	\end{pmatrix}
	                                                    & \\
	            & \Longleftrightarrow
	\begin{pmatrix}
	  a & -8 c \\
	  b & -8 d
	\end{pmatrix}
	=
	\begin{pmatrix}
	  d  & -b \\
	  8c & -8a
	\end{pmatrix}
	                                                    & \\
	            & \Longleftrightarrow a = d \text{ et } b = 8c
\end{flalign*}


\medskip

La condition $\det A = 1$ pour
$A
=
\begin{pmatrix}
  a & b \\
  c & d
\end{pmatrix}
=
\begin{pmatrix}
  a & 8c \\
  c & a
\end{pmatrix}$
nous donne
$a^2 - 8c^2 = 1$,
autrement dit, $(x;y) = (a:c)$ doit être une solution de \ED. Que c'est joli !


\medskip

Notons que l'ensemble des matrices de ce type est stable par multiplication,
\footnote{
    La démonstration est faite un peu plus bas dans ce document.
}
et que la matrice du sujet de Bac n'est autre que celle correspondant à la solution élémentaire
$(a \,; c) = (3 \,; 1)$.%
\footnote{
    Le choix de $(a \,; c) = (1 \,; 0)$ nous conduit à la matrice identité qui ne nous fournira aucune nouvelle solution à partir d'une existante.
}


% ------------------ %


\medskip

\textit{\textbf{Cas 2:} supposons maintenant $\det A = -1$.}

\medskip

Notant
$A = \begin{pmatrix}
  a & b \\
  c & d
\end{pmatrix}$,
nous avons alors
$A^{-1} = \begin{pmatrix}
  -d & b  \\
   c & -a
\end{pmatrix}$,
d'où comme dans les calculs précédents:
$A^T Q A = Q \Longleftrightarrow a = -d \text{ et } b = -8c$.


\medskip

La condition $\det A = -1$ pour
$A
=
\begin{pmatrix}
  a & b \\
  c & d
\end{pmatrix}
=
\begin{pmatrix}
  a & -8c \\
  c & -a
\end{pmatrix}$
nous redonne
$a^2 - 8c^2 = 1$,
mais avec un autre ensemble de matrices qui lui n'est pas stable par multiplication.%
\footnote{
    Ici, c'est immédiat, car $\det A = -1$ et $\det B = -1$ donnent 
    $\det(AB) = \det A \det B = 1$.
}
